\section{Nine complete small orthogonal core reductions}
\label{sec:orthogonal-small-cores}

\paragraph{Status and scope.}
New root derivation. Complete author and independent peer proof and
source reviews have passed, together with separate full exact spectrum
reconstructions and the retained GAP orthogonality checks. The character tables are
cited exact inputs, not independently reconstructed original tables.
Eight simple types are excluded in arbitrary mixed products. The ninth,
$\Omega_9(3)$, has a separately stated transitive-core conclusion in the
nonselfdual-original inducing problem. No assertion about arbitrary
mixtures involving that ninth type is substituted for its hypotheses.

Write $\nu(T)=a-\max_\lambda\dim\ker(T-\lambda I)$ for
projective support on a complex space of degree $a$.

\begin{lemma}[D-OSM-01: complete ordinary spectra]
\label{D-OSM-01}
For each complete ordinary table in the following display, take every
nonlinear irreducible row. For an extension table take precisely its
rows whose restriction to the displayed universal core is irreducible.
The minimum degree and the minimum $\nu/a$ over all nonscalar elements
are exactly as follows.
\[
\begin{array}{l|r|r|r|c}
\text{table}&\text{rows}&\text{spectra}&\min a&\min\nu/a\\\hline
6.O_7(3)&179&31146&27&2/9\\
6.O_7(3).2&126&22932&78&11/39\\
O_8^-(2)&38&1444&34&16/51\\
O_8^-(2).2&52&3068&34&7/34\\
2^2.O_8^+(3)&353&123550&260&28/65\\
2.O_8^-(3)&189&35532&246&52/123\\
2.O_9(3)&325&105300&780&13/41\\
O_{10}^+(2)&96&9216&155&32/93\\
O_{10}^+(2).2&140&21560&155&7/31\\
O_{10}^-(2)&114&12996&154&64/187\\
O_{10}^-(2).2&140&22820&154&5/22\\
O_{12}^+(2)&270&72900&651&128/357\\
O_{12}^+(2).2&340&132940&651&5/21\\
O_{12}^-(2)&236&55696&650&256/715\\
O_{12}^-(2).2&340&127160&650&31/130
\end{array}                                                     \tag{OSM.1}
\]
The only nonselfdual rows of $2.O_9(3)$ are rows $224,225$,
both of degree $838656$ and minimum support ratio $1/2$.
All its other rows have Frobenius--Schur indicator $+1$.
\end{lemma}

\begin{proof}
The mathematical inputs are the complete official CTblLib tables
\cite{HD:OrthogonalSmallTables}. The exporter retains every character
value, class size and order, every prime power map, all six core
fusions, all indicators, and the original provenance strings.
The complete input has SHA-256
\texttt{5236207618c2564af7ed4871c0f92a667664255d2e60be87f037257356669b79}.
GAP verifies full row orthogonality for all fifteen tables. A GAP
spectrum computation and a separate exact rational-trace program
recover every eigenvalue multiplicity by
\[
 m_k(x)=\frac1d\sum_{j=0}^{d-1}\chi(x^j)\zeta_d^{-kj},
 \qquad d=|x|.                                                \tag{OSM.2}
\]
The separate program selects extension rows by restriction inner
product one and checks exact pointwise equality with a unique core
row. It checks full class unions, the identity kernel, element orders,
and all retained power-map compatibilities for the fusions. All
Fourier multiplicities are exact nonnegative integers with the
correct degree sum. Scalar elements are detected by
$|\chi(x)|^2=\chi(1)^2$; no inertia class is omitted.

There are exactly $2938$ selected rows and $778260$ nonscalar spectra.
The retained per-row witnesses prove all the displayed minima.
The separate program also recomputes every Frobenius--Schur indicator
using $|G|^{-1}\sum_x\chi(x^2)$, with exact rational traces of the
cyclotomic values. The original sum is a rational integer, so taking
its normalized rational trace preserves it. This verifies the two
exceptional nonselfdual rows and the indicators of every other row.

The strings $O_d^\epsilon(q)$ here denote the simple groups in the
official table notation, not the full natural orthogonal groups.
The retained multiplier/outer metadata, interpreted as in
\cite[\S3.7-3]{HD:OrthogonalSmallTables}, are respectively
\[
\begin{array}{c|c|c}
S&M(S)&\operatorname{Out}(S)\\\hline
O_7(3)&6&2\\
O_8^-(2)&1&2\\
O_8^+(3)&2^2&S_4\\
O_8^-(3)&2&2^2\\
O_9(3)&2&2\\
O_{10}^{\pm}(2)&1&2\\
O_{12}^{\pm}(2)&1&2
\end{array}                                                     \tag{OSM.3}
\]
The exceptional $\Omega_7(3)$ multiplier is also located in
\cite[Chapter~5, Table~1 and preceding paragraph, p.~35]{BI:HHM}.
Every quasisimple central quotient lifts to the stated universal cover.
Keeping all nonlinear rows therefore includes all central characters,
including their nonfaithful central kernels. Faithfulness on the
noncyclic center of $\operatorname{Spin}_8^+(3)$ is not assumed.
\end{proof}

\begin{lemma}[D-OSM-02: full normalizers]
\label{D-OSM-02}
For every quasisimple core with simple quotient among
\[
\mathcal F=\{\Omega_7(3),\Omega_8^-(2),
 \mathrm{P}\Omega_8^+(3),\mathrm{P}\Omega_8^-(3),
 \Omega_{10}^+(2),\Omega_{10}^-(2),
 \Omega_{12}^+(2),\Omega_{12}^-(2)\},                         \tag{OSM.4}
\]
every finite-order nonscalar normalizer of a nontrivial irreducible
degree-$a$ representation has support at least $a/5$.
For a core with simple quotient $\Omega_9(3)$ the bound is
$13a/82$; it improves to $a/4$ when the core representation is
nonselfdual.
\end{lemma}

\begin{proof}
First consider the six types in (OSM.4) other than the two
eight-dimensional natural types over $\mathbf F_3$.
Their outer groups have order two, and the complete universal core
and covering outer table are both in (OSM.1). An automorphism of a
represented central quotient lifts to the universal perfect cover.
Such a lift is unique: an automorphism trivial on the simple quotient
differs from the identity by a homomorphism of a perfect group to its
center, which is trivial.

Let $T$ implement a compatible outer automorphism. Choose an element
$t$ in the retained covering extension inducing the same automorphism,
after multiplication by a core element. By Schur's lemma
$T^2=c\rho(t^2)$ for a nonzero scalar $c$. Rescaling $T$ solves this
actual square relation, so $\rho$ extends to a representation of the
listed covering extension with that action. It is irreducible because
its restriction is irreducible. Every such row is retained. This
argument permits a nonsplit relation and an isoclinic alternative;
it does not assert that isoclinic groups have equal character tables.
Changing scalar phases preserves projective support.
The minimum across these six complete core/outer pairs exceeds $1/5$.

For the two remaining types in (OSM.4), every nonscalar core operator
has support at least $(2/5)a$, since $28/65,52/123>2/5$.
A nonscalar normalizer $T$ fails to commute with some represented
core element $g$, by irreducibility. If $[T,g]$ is nonscalar its
actual moving rank is at least $(2/5)a$. If it is a nonidentity
scalar its actual moving rank is $a$. Scalar-normalize $T$ at its
largest eigenspace. Then
\[
 \operatorname{rank}([T,g]-I)
 =\operatorname{rank}(Tg-gT)\le2\operatorname{rank}(T-I)
 =2\nu(T).                                                    \tag{OSM.5}
\]
This proves $\nu(T)\ge a/5$ for every compatible normalizer,
including all outer automorphisms of the noncyclic central cover.
No table for the full $S_4$ outer group is required.

The identical commutator argument applied to $2.O_9(3)$ gives
$\nu(T)\ge13a/82$. For its two nonselfdual characters, the full
core support is $a/2$, so the argument gives $\nu(T)\ge a/4$.
Every central quotient is already covered by the universal rows.
\end{proof}

\begin{theorem}[D-OSM-03: eight mixed-type source exclusions]
\label{D-OSM-03}
An irreducibly acting normal central-product core whose simple
quotients all lie in $\mathcal F$ cannot supply actual period rank
nine in the signature-$(3,3)$ sixfold-source problem. Arbitrary
mixtures of those eight types, covers, scalars, compatible outer
operators and factor permutations are allowed. There is no
selfduality or determinant hypothesis.
\end{theorem}

\begin{proof}
The core module is a tensor product of nontrivial irreducible factor
modules. A factor-preserving nonscalar operator has normalized
eigenspace multiplicities given by convolution; their maximum is at
most the maximum for any nonscalar factor. D-OSM-02 therefore gives
support at least $e/5$, where $e$ is the total degree.
For a nontrivial factor cycle of length $r$ and factor degree $a$,
conjugate the intertwiners to identities except for a diagonal closing
twist. This is a weighted cyclic permutation on $a^r$ basis words.
There are $a$ constant words, and every other orbit has length at
least two. Each orbit contributes at most one to any chosen
eigenspace, so its dimension is at most
$a+(a^r-a)/2\le3a^r/4$. Tensoring other cycles cannot increase
the normalized maximum. The resulting quarter support is stronger
than required. A genuine nonidentity scalar inertia moves all $e$
dimensions.

Every factor has degree at least $27$. The actual intermediate-cover
identity, with base genus $h$ and $R$ genuine inertia points, is
$6=2e(h-1)+D$, where $D$ is the sum of actual moving ranks.
Genus $h\ge2$ is impossible; genus one has $D=6$, hence at most six
effective branches and actual period-image dimension at most six.
Rank nine would force genus zero, $R\ge12$ and $D=2e+6$.
Consequently $2e+6\ge12e/5$, giving $e\le15<27$, a contradiction.
This uses the actual variation factorization, valid for a larger
inducing character field; it does not infer dominance from a parameter
count.
\end{proof}

\begin{theorem}[D-OSM-04: the ninth transitive core]
\label{D-OSM-04}
Suppose the original sixfold-source representation is nonselfdual,
and an inducing block has a transitive irreducibly acting normal
central-product core with simple factors $\Omega_9(3)$. It cannot
have an actual period image of dimension nine. Every compatible
normalizer, scalar, factor permutation and either determinant is allowed.
\end{theorem}

\begin{proof}
Transitivity makes all core factors either selfdual or nonselfdual.
If they are nonselfdual, D-OSM-02 and the tensor-cycle argument give
quarter support for every nonscalar inertia. Rank nine would force
$2e+6\ge3e$, hence $e\le6$, whereas every factor has degree at
least $780$.

Otherwise each factor has an invariant nondegenerate bilinear form,
unique up to scalar. Their tensor product gives the unique invariant
form line for the irreducible core. The inducing image acts on it by
a multiplier character $\mu$. This character is nontrivial: if it
were trivial, the inducing representation would be selfdual, and
its finite-group induction would make the original representation
selfdual, contrary to hypothesis. No equality of character fields is
used.

The same actual Euler/Hurwitz argument forces genus zero and at least
twelve genuine inertias, generating the inducing image and having
product one. At least two have nontrivial $\mu$. Their fixed spaces
are totally isotropic for the form and their moving ranks are at least
$e/2$. Every remaining genuine inertia costs at least $13e/82$ by
D-OSM-02 and the tensor-cycle bound. Thus
\[
 2e+6\ge e+10(13e/82),\qquad e\le41/4<780,                 \tag{OSM.6}
\]
again impossible. The selfdual/nonselfdual division is essential to
this argument; arbitrary intransitive mixed factors involving this
ninth type are not asserted to be excluded.
\end{proof}

\paragraph{Boundary and geometric placement.}
These are necessary restrictions on finite-cover weight-one variations,
their integral representation-ring operations and actual period maps.
They narrow the condition-A cycle-supplier gate at cote16, scan050,
in the EGA/SGA and realization framework stated in the paper. They
produce no cycle, positivity division, Cayley--Hamilton projector or
motivic $L$-value estimate. Other fields and orthogonal ranks, in
particular the small $\Omega_8^+(2)$ cover representations, retain
their separate gates. No Markman, Severn or 032 premise is used.
